\chapter*{Where To Go Next}
\markboth{Where To Go Next}{Where To Go Next}
\phantomsection
\addcontentsline{toc}{chapter}{Where To Go Next}

When you finish a program, you are not really finished. A working program is a
place to start experimenting.

\section*{Five Good Changes}
\phantomsection
\addcontentsline{toc}{section}{Five Good Changes}

\begin{enumerate}
\item Change the colors.
\item Add a sound when something good happens.
\item Add a sound when something bad happens.
\item Make the game a little faster after each point.
\item Add a title screen line with your name on it.
\end{enumerate}

\section*{Bigger Challenges}
\phantomsection
\addcontentsline{toc}{section}{Bigger Challenges}

After the small changes, try one bigger change:

\begin{itemize}
\item Add a second player.
\item Add a high score.
\item Add a new enemy pattern.
\item Draw a better board or background.
\item Save a level or score table to disk.
\end{itemize}

\section*{Keep A Programmer's Notebook}
\phantomsection
\addcontentsline{toc}{section}{Keep A Programmer's Notebook}

Write down the changes you try. Keep the ones that work. Cross out the ones
that do not. A notebook turns random experiments into a map of what you have
learned.

\begin{typeinbox}[title=One Last Secret]
Nobody understands a whole program all at once. Good programmers understand one
small piece, then another, then another. That is enough.
\end{typeinbox}
